\section{Basis Set Comparison}\label{sec:comparison}
% ============================================================
% ============================================================

\begin{center}
\renewcommand{\arraystretch}{1.5}
\begin{tabular}{lcccc}
  \toprule
  Property & Basis I (HO) & Basis II (Shifted HO) & Basis III (Gaussian) & Basis IV (DVR) \\
  \midrule
  Orthogonal          & Yes & No  & No  & Yes \\
  Overlap $\mathbf{S}$ & $\mathbf{I}$ & Non-trivial & Non-trivial & $\mathbf{I}$ \\
  $\mathbf{H}$ structure & Banded & Dense & Dense & Dense \\
  Parity factorization & Yes (symmetric $V$) & Possible & No & No \\
  $\mathbf{V}$ matrix  & Analytic & Analytic & Analytic & Diagonal \\
  Parameters           & $\alpha$ & $\alpha, \xe$ & $\{x_i, \alpha_i\}$ & $\Delta x, N$ \\
  Typical $N$ for $10^{-8}$ & 30--50 & 15--25 & 20--40 & 200--500 \\
  Benchmark quality    & No & No & No & Yes \\
  \bottomrule
\end{tabular}
\end{center}

% ------------------------------------------------------------
\subsection{Analytical comparison for the asymmetric double well}\label{sec:basis_asym}
% ------------------------------------------------------------

The summary above compares the four bases in generic terms.  For the \emph{asymmetric}
potential $V(x) = ax^4 - bx^2 + cx$ ($c \neq 0$), several properties change qualitatively,
making the DVR (Basis~IV) analytically preferable to the shifted-HO (Basis~II) despite
Basis~II achieving faster convergence per function.  This subsection examines that
trade-off in detail.

\subsubsection{Matrix element complexity}

\textbf{DVR/PIB basis.}
The kinetic energy matrix is given by the closed-form Colbert--Miller expressions
(\cref{eq:T_DVR_diag,eq:T_DVR_offdiag}).  The potential matrix is diagonal:
$V_{jl} = V(x_j)\,\delta_{jl}$ (\cref{eq:V_DVR}).  No integral is ever computed;
the Hamiltonian is assembled solely from point evaluations of $V(x)$ and the analytic
kinetic-energy formula.  Changing the potential requires changing only the function
evaluation, not any integral routine.

\textbf{Shifted-HO basis.}
Same-center matrix elements reduce to origin-centered HO results via the binomial shift
(\cref{sec:shifted_potential}), but the cross-center blocks $\langle\varphi_m^{(L)}|x^k|
\varphi_n^{(R)}\rangle$ require: (i)~computing the exponential overlap factor
$S_{mn}^{LR}$ involving associated Laguerre polynomials (\cref{eq:S_LR}); (ii)~expanding
$x^k$ in displacements about the right center; and (iii)~contracting the result with
$S_{mn}^{LR}$.  The bookkeeping grows with $k$ and becomes substantially heavier when
$\alpha_L \neq \alpha_R$ (which is the generic asymmetric case, see below).

\subsubsection{Orthogonality and the eigenvalue problem}

The DVR basis is orthonormal, giving the standard eigenvalue problem $\mathbf{H}\mathbf{c}
= E\mathbf{c}$.

The shifted-HO basis is non-orthogonal, requiring the generalized problem
$\mathbf{H}\mathbf{c} = E\,\mathbf{S}\mathbf{c}$.  The overlap matrix $\mathbf{S}$ must
be computed, stored, and factored (e.g.\ via Cholesky or L\"owdin orthogonalization
$\mathbf{S}^{-1/2}$) before the standard solver can be applied.  This introduces an
additional source of numerical ill-conditioning: if the two well centers are close (shallow
barrier, strong asymmetry pushing $x_-$ and $x_+$ together), the cross-center overlaps
$S_{mn}^{LR}$ approach unity and $\mathbf{S}$ develops near-linear dependencies.  A
regularization or basis truncation is then required.

\begin{remark}
  The variational theorem (Hylleraas--Undheim--MacDonald) holds for the generalized
  eigenvalue problem, so monotone convergence with basis size is guaranteed in principle.
  In practice, however, near-linear dependence limits the useful basis size before
  numerical noise dominates.
\end{remark}

\subsubsection{Asymmetry-specific complications for Basis~II}

In the \emph{symmetric} case the shifted-HO basis benefits from:
\begin{itemize}
  \item $x_L = -\xe$, $x_R = +\xe$ known analytically;
  \item $\alpha_L = \alpha_R = \alpha_{\text{local}}$ by symmetry;
  \item a parity decomposition that halves the dimension of each block.
\end{itemize}

For $c \neq 0$, all three advantages are lost:
\begin{enumerate}
  \item \textbf{Shifted centers.}  The minima $x_\pm$ are roots of the depressed cubic
    $4ax^3 - 2bx + c = 0$ (\cref{eq:cubic_roots}), which must be solved (numerically
    or via the trigonometric formula) at each value of $c$.
  \item \textbf{Unequal widths.}  The curvatures $V''(x_\pm) = 12ax_\pm^2 - 2b$ differ,
    so $\omega_L \neq \omega_R$ and the natural width parameters $\alpha_L = \mu\omega_L$,
    $\alpha_R = \mu\omega_R$ must be computed separately.  With $\alpha_L \neq \alpha_R$,
    the displacement-operator formula \cref{eq:S_LR} no longer applies directly; the
    overlap integrals must be evaluated by combining two Gaussians of different widths,
    leading to a more involved computation.
  \item \textbf{No parity block structure.}  The linear term $cx$ couples even and odd
    functions (\cref{sec:parity}), so the full $N\times N$ generalized eigenvalue problem
    must be solved without any block decomposition.
  \item \textbf{Basis partitioning.}  The optimal split $N_L$ vs.\ $N_R$ between the two
    centers is not obvious a priori when the wells are unequal, and an inappropriate
    partition may under-represent one well.
  \item \textbf{Discontinuous basis under parameter variation.}  A sweep over $c$ requires
    recomputing centers, widths, and overlaps at every parameter value, making the basis
    non-uniform across the sweep.
\end{enumerate}

The DVR basis is entirely unaffected by all five points: the same grid with the same
parameters applies for any $c$.

\subsubsection{The DVR/FGH connection: elimination of all integrals}

The deepest analytical advantage of the DVR basis is the
\emph{integral-free representation} of any potential.  The Fourier Grid Hamiltonian
(\cref{sec:FGH}) exploits the fact that the PIB/sinusoidal functions define a unitary
transformation between the finite basis representation (FBR) and a uniform grid of
quadrature points $\{x_j\}$.  In the DVR, the potential matrix becomes exactly diagonal
(\cref{eq:V_DVR}):
\begin{equation}
  V_{jl} = V(x_j)\,\delta_{jl}\,.
\end{equation}
No integrals need to be computed for any potential, polynomial or otherwise.  The
Hamiltonian matrix is assembled in $\mathcal{O}(N^2)$ operations (one potential
evaluation per grid point, plus the analytic kinetic-energy formula).

The shifted-HO basis has no natural DVR.  A Gauss--Hermite quadrature exists for
each individual center, but combining two centers with different widths requires a
non-standard bi-centric quadrature that negates the simplicity of the approach.

\subsubsection{Time-dependent dynamics}

For wavepacket propagation studies (e.g.\ survival probability, expectation values as
functions of time), the DVR/PIB basis connects naturally to the \emph{split-operator}
method via the Fast Fourier Transform:
\begin{equation}
  e^{-i\hat{H}\Delta t/\hb}
  \approx e^{-i\hat{V}\Delta t/(2\hb)}\,\mathcal{F}^{-1}\,
          e^{-i\hat{T}\Delta t/\hb}\,\mathcal{F}\,
          e^{-i\hat{V}\Delta t/(2\hb)}
  + \mathcal{O}(\Delta t^3)\,,
\end{equation}
where $\mathcal{F}$ is the discrete Fourier transform.  Each step costs
$\mathcal{O}(N\log N)$.  The potential exponential is diagonal in the position
representation; the kinetic exponential is diagonal in momentum space.  This requires
no matrix diagonalization whatsoever.

For the shifted-HO basis, the propagator $e^{-i\mathbf{S}^{-1}\mathbf{H}\Delta t/\hb}$
must be computed as a dense matrix exponential, costing $\mathcal{O}(N^3)$ per
time step.

\subsubsection{Summary}

\begin{center}
\renewcommand{\arraystretch}{1.4}
\begin{tabular}{lcc}
  \toprule
  Criterion & DVR (Basis~IV) & Shifted-HO (Basis~II) \\
  \midrule
  Convergence rate (eigenstates needed) & Slower & \textbf{Faster} \\
  Matrix element simplicity & \textbf{Integral-free} & Heavy (cross-center) \\
  Eigenvalue problem type & \textbf{Standard} & Generalized \\
  Overlap conditioning & \textbf{None} & Risk near $\Delta E \to 0$ \\
  Parity exploitation ($c=0$) & No & Yes \\
  Robustness for $c \neq 0$ & \textbf{Full} & Loses parity; re-parameterize \\
  Parameter sweep ($c$ varies) & \textbf{Trivial} & Recompute basis each step \\
  Time-dependent propagation & \textbf{FFT split-op} $\mathcal{O}(N\log N)$ & Dense expm $\mathcal{O}(N^3)$ \\
  Physical transparency & Less & \textbf{More} \\
  \bottomrule
\end{tabular}
\end{center}

\textbf{Conclusion.}
The shifted-HO basis (Basis~II) is physically transparent and achieves excellent
convergence per function in the symmetric case, where the well centers and widths are
known analytically and parity halves the problem size.  In the asymmetric case these
advantages largely disappear: the basis must be redesigned at every parameter value,
the generalized eigenvalue problem introduces conditioning risk, and the integral-free
simplicity of the DVR is foregone.

The DVR/FGH basis (Basis~IV) requires modestly more grid points to reach the same
accuracy, but this cost is offset by: (i)~zero integral computation for any potential;
(ii)~a standard, well-conditioned eigenvalue problem; (iii)~complete insensitivity to
$c$; and (iv)~the ability to perform time-dependent propagation at $\mathcal{O}(N\log N)$
cost.  For the asymmetric double well---and particularly for parameter sweeps over the
asymmetry coefficient $c$---the DVR is the analytically and computationally preferred basis.

% ============================================================
% ============================================================
